% Einheit 3 von 5 – Prozentwert berechnen (bank/prozentrechnung/e3.jsonl, zusammenbau v0.1)
%% TODO Zweigzeile Teil 1: was hier gelernt wird (Satz in Schülersprache aus dem Einheitstitel) – Einheit 3
\einheitenkopf[e3]{Einheit 3 von 5 · Prozentwert berechnen}
\zweigzeile{ab Kl. 7 · P10}
\setcounter{aufgabe}{16}

%% TODO Titel: Ich-kann-Satz fehlt in der Bank (Platzhalter: Kettenname) – Nr. 17
%% TODO Anweisung: ein Satz über der ersten Teilaufgabe fehlt in der Bank – Nr. 17
\begin{aufgabe}{Prozentwert}
\swa{$50\,\%$ von $80\,€$ – wie viel? \leerfeld[€]}{\streifen[2]{50}{$0$}{$80\,€$}}
\swa{$25\,\%$ von $60$ kg – wie viel? \leerfeld[kg]}{\streifen[4]{25}{$0$}{$60$ kg}}
\swa{$10\,\%$ von $300$ m – wie viel? \leerfeld[m]}{\streifen{10}{$0$}{$300$ m}}
\swa{$50\,\%$ von $14$ l – wie viel? \leerfeld[l]}{\streifen[2]{50}{$0$}{$14$ l}}
\swa{$10\,\%$ von $90\,€$ – wie viel? \leerfeld[€]}{\streifen{10}{$0$}{$90\,€$}}
\swfrage{Was bleibt gleich, was ändert sich?}
%% TODO Darstellung für schwach fehlt in der Bank (prozentrechnung-e3-k1-s2-v1; Abschnitt „Für schwache Schüler“)
\swz{$6\,\%$ von $400$ kg}{}
%% TODO Darstellung für schwach fehlt in der Bank (prozentrechnung-e3-k1-s3-v1; Abschnitt „Für schwache Schüler“)
\swz{$30\,\%$ von $90\,€$}{}
%% TODO Darstellung für schwach fehlt in der Bank (prozentrechnung-e3-k1-s4-v1; Abschnitt „Für schwache Schüler“)
\swz{$45\,\%$ von $80$ kg – mit Dezimalzahl}{}
%% TODO Darstellung für schwach fehlt in der Bank (prozentrechnung-e3-k1-s5-v1; Abschnitt „Für schwache Schüler“)
\swz{$20\,\%$ von $4{,}50\,€$}{}
%% TODO Darstellung für schwach fehlt in der Bank (prozentrechnung-e3-k1-s6-v1; Abschnitt „Für schwache Schüler“)
\swz{Schuhe für $70\,€$, $20\,\%$ Rabatt – wie viel spart man?}{}
%% TODO Darstellung für schwach fehlt in der Bank (prozentrechnung-e3-k1-s7-v1; Abschnitt „Für schwache Schüler“)
\swz{$150\,\%$ von $40$ kg}{}
\begin{teile}
\teil Ein Fernseher kostet $480\,€$. Bei Barzahlung gibt es $15\,\%$ Rabatt. Wie viel Euro spart man? Kreuze an. \hfill (P10 2021 OS)\\ \kreuz{$32\,€$}\\ \kreuz{$408\,€$}\\ \kreuz{$72\,€$}\\ \kreuz{$552\,€$}
\end{teile}
\swz[3]{$17\,\%$ von $60\,€$ \hfill (P10 2014 OS)}{}
\swz[3]{Bei einer Umfrage unter $800$ Haushalten hatten $35\,\%$ ein E-Bike, sechs Jahre später $52\,\%$. Um wie viele Haushalte hat die Zahl zugenommen? \hfill (P10 2019 OS)}{}
\end{aufgabe}

%% TODO Titel: Ich-kann-Satz fehlt in der Bank (Platzhalter: Kettenname) – Nr. 18
%% TODO Anweisung: ein Satz über der ersten Teilaufgabe fehlt in der Bank – Nr. 18
\begin{aufgabe}{Mehrwertsteuer in Euro}
%% TODO Darstellung für schwach fehlt in der Bank (prozentrechnung-e3-k2-s1-v1; Abschnitt „Für schwache Schüler“)
\swz{Reparatur ohne Mehrwertsteuer $300\,€$, dazu $19\,\%$ – wie viel Euro Steuer?}{}
\end{aufgabe}

%% TODO Titel: Ich-kann-Satz fehlt in der Bank (Platzhalter: Kettenname) – Nr. 19
%% TODO Anweisung: ein Satz über der ersten Teilaufgabe fehlt in der Bank – Nr. 19
\begin{aufgabe}{Prozentwert – Fehler finden, Begründen, Darstellungswechsel, Anwendung}
\swz[3]{Jonas, $60\,\%$ von $35\,€$: \rechnung{0{,}06 \cdot 35 &= 2{,}10\,€} Fehler? Wie heißt es richtig?}{}
\swfrage{$1\,\%$ von $500\,€$ sind $5\,€$. Warum?}
\swa{Ganzes $60\,€$: $30\,\%$ im Streifen einzeichnen, Wert ablesen. \leerfeld[€]}{\streifen[0]{0}{$0$}{$60\,€$}}
\swz[3]{Dieselbe Jacke: in Laden A $95\,€$ mit $20\,\%$ Rabatt, in Laden B $85\,€$ mit $10\,\%$ Rabatt. Wo ist sie billiger?}{}
\end{aufgabe}

\uebersichtskasten{Prozentwert: erst 1 \% (Ganzes : 100), dann mal p – oder p als Dezimalzahl mal Ganzes. \\ 18 \% von 40 kg: 1 \% = 0,4 kg → 18 \% = 7,2 kg \quad 0,18 · 40 = 7,2}
